\chapter{Scheduling}
\label{ch:scheduling}

Whenever more threads are ready than there are CPUs, the scheduler
decides who runs --- at every timer interrupt, and whenever a thread
blocks, wakes up, yields or exits. This chapter covers the metrics, the
classic policies you will be asked to simulate by hand in the exam, how
real systems combine them, and the question that matters most for your
\code{pthread\_join}: should a waiting thread spin, yield, or sleep?

% \gantt{chart}: one box per character, '.' = idle
\newcommand{\gantt}[1]{%
\begin{tikzpicture}[x=0.52cm,y=0.5cm,font=\scriptsize]
  \def\t{0}%
  \foreach \c [count=\i from 0] in {#1} {%
    \ifx\c.\draw (\i,0) rectangle ++(1,1);%
    \else\draw[fill=newbg] (\i,0) rectangle ++(1,1) node[pos=0.5] {\c};\fi
    \node[below] at (\i,0) {\i};
  }%
\end{tikzpicture}}

\section{Bursts and metrics}

Threads alternate between \emph{CPU bursts} (computing) and \emph{I/O
bursts} (waiting for the disk, the network, the keyboard, a lock).
CPU-bound threads have long CPU bursts, I/O-bound (interactive) threads
short ones. For a job that arrives at time $a$, first runs at $s$ and
finishes at $f$ with total CPU time $b$:

\begin{center}
\begin{tabular}{lll}
\toprule
\textbf{Metric} & \textbf{Definition} & \textbf{Who cares} \\
\midrule
turnaround time & $f - a$ & batch jobs \\
waiting time & $f - a - b$ (time spent ready, not running) & fairness \\
response time & $s - a$ & interactive users \\
throughput & jobs finished per time unit & the whole system \\
\bottomrule
\end{tabular}
\end{center}

The goals conflict: short quanta improve response and cost throughput
(context switches); favouring short jobs improves the average and can
starve long ones.

\section{The classic policies}

Example (arrival, burst, priority): A (0, 5, 2), B (1, 3, 1), C (2, 1, 3),
D (4, 2, 1) --- \module{09}, question 1. One box per time unit.

\paragraph{FCFS} (first come, first served): run jobs to completion in
arrival order. Simple, no starvation, but short jobs wait behind long ones
--- the \emph{convoy effect}.

\begin{center}\gantt{A,A,A,A,A,B,B,B,C,D,D}\\ average waiting 3.75\end{center}

\paragraph{SJF} (shortest job first, non-preemptive): when the CPU becomes
free, pick the job with the shortest burst. When all jobs are present
from the start, this is provably optimal for the average waiting time
(exchange argument: a longer job directly in front of a shorter one ---
swap them, and the sum of waiting times drops). With staggered arrivals
it is only a good heuristic: a short job arriving just after a long one
has started still has to wait.

\begin{center}\gantt{A,A,A,A,A,C,D,D,B,B,B}\\ average waiting 3.00\end{center}

\paragraph{SRTF} (shortest remaining time first): SJF with preemption ---
a new job with less remaining work than the running one takes over.
Optimal for the average waiting time over all schedules.

\begin{center}\gantt{A,B,C,B,B,D,D,A,A,A,A}\\ average waiting 2.00, response 0.25\end{center}

Both need the length of the next CPU burst, which the OS does not know.
It can \emph{predict} it by exponential averaging of the past bursts:
$\tau_{n+1} = \alpha t_n + (1-\alpha)\tau_n$.

\paragraph{Round robin} (RR): a FIFO ready queue; each job runs for at
most one \emph{quantum} $q$, then goes to the end of the queue. Good
response time, fair, no starvation. The quantum is the knob: too small and
the CPU is busy switching (a 0.1~ms switch with $q$ = 1~ms wastes 9\%),
too large and RR degenerates into FCFS. Typical values: a few
milliseconds. With $q = 2$:

\begin{center}\gantt{A,A,B,B,C,A,A,D,D,B,A}\\ average waiting 4.25, response 1.50\end{center}

A tie that exams love: at $t = 2$, C arrives \emph{and} A's quantum
ends. The convention used in \module{09} (and by most textbooks): the
newcomer is queued first, the preempted job behind it. State your
convention when you answer.

\paragraph{Priority} scheduling: always run the most important ready job
(here: smaller number = more important; preemptive):

\begin{center}\gantt{A,B,B,B,D,D,A,A,A,A,C}\\ average waiting 3.25\end{center}

Low-priority jobs can \emph{starve} if higher-priority work keeps
arriving. Remedy: \emph{aging} --- raise the priority of waiting jobs over
time. Also watch for priority inversion (chapter~6).

\section{Multi-level feedback queue}

MLFQ learns what kind of job it deals with, without being told:

\begin{enumerate}
  \item Several queues with decreasing priority; round robin within a
    queue; higher queues always win, and preempt lower ones.
  \item A new job enters the top queue.
  \item A job that uses up its time allotment on a level is moved one
    level down. The allotment is counted \emph{across} I/O --- otherwise a
    job could stay on top by blocking just before its quantum ends.
  \item Periodically, all jobs are boosted back to the top: no
    starvation, and a job that changed its behaviour gets a new chance.
\end{enumerate}

Interactive jobs (short bursts) stay high and get quick responses;
CPU-bound jobs sink and share the rest. \module{09}'s bonus implements it
with exact rules.

\section{Proportional share and Linux}

Instead of priorities, give each job a \emph{share}: lottery scheduling
draws tickets, stride scheduling advances a ``pass'' value by
1/tickets. Linux's scheduler (CFS, since 6.6 EEVDF) is of this kind: each
task accumulates \emph{virtual runtime}, weighted by its nice value (nice
0 = weight 1024, each nice step $\approx$ 1.25$\times$ less), and the task
with the least virtual runtime runs next. A task that sleeps a lot has
little virtual runtime and runs soon after waking --- interactive tasks
stay responsive without explicit classification (\module{09}'s examples
measure both effects).

\section{Spin, yield or sleep?}

A thread that must wait for another thread (a flag, a lock, a join) on a
\emph{single} CPU:

\begin{description}
  \item[spin] (\code{while (!flag) \{\}}): burns its whole quantum every
    round --- and the flag cannot change meanwhile, because the thread
    that would set it is not running. Pure waste.
  \item[yield] (\code{while (!flag) yield();}): gives the CPU away at once,
    but stays \emph{runnable}: it is scheduled again and again just to
    check and yield. Costs a context switch per round, forever.
  \item[sleep] (block until woken): leaves the ready queue entirely; costs
    exactly one switch away and one back.
\end{description}

A kernel waits by sleeping, except for very short waits on locks held
for a few instructions.

\begin{swebbox}
\code{Scheduler::schedule()} (\file{common/source/kernel/Scheduler.cpp})
is round robin: it picks the first thread in \code{threads\_} that is
\code{Running} (i.e.\ schedulable) and rotates the list. The timer
interrupt (\code{irqHandler\_0}) calls it every tick, so the quantum is one
tick; \code{Scheduler::yield()} triggers the same code voluntarily with
\code{int \$65}. When nothing is runnable, the idle thread runs.
\code{sleep()} sets the state to \code{Sleeping}, which removes a thread
from consideration until \code{wake()} sets it back. For
\code{pthread\_join}: the joiner must \emph{sleep} until the target has
exited --- a \code{while (!done) yield();} loop would ``work'' and waste
every round-robin cycle.
\end{swebbox}

\begin{keybox}
\begin{itemize}
  \item Metrics: turnaround $f-a$, waiting $f-a-b$, response $s-a$.
  \item FCFS (convoys), SJF (optimal when all jobs start together), SRTF (optimal),
    RR (response; quantum trade-off), priority (starvation $\to$ aging).
  \item MLFQ: demote on used allotment (across I/O), boost periodically.
  \item Linux: weighted virtual runtime, nice changes the weight.
  \item Wait by sleeping; on one CPU, spinning is pointless and yielding
    in a loop is wasteful.
\end{itemize}
\end{keybox}
